\subsection{Same-state reconstruction and explicit rephasing}
\label{sec:quadratic_same_state_converse}

\begin{lemma}[Rephasing within the same quadratic system]
    \label{lem:quadratic_solution_rephasing}
    \lean{MPSTensor.QuadraticReconstructionSolution.rephase,
        MPSTensor.QuadraticReconstructionSolution.rephase_tensor,
        MPSTensor.QuadraticReconstructionSolution.mpv_rephase}
    \notready
    \uses{def:pgvwc07_quadratic_system}
    Let $(B,Y_j,Z_j)$ solve the quadratic system and let $|\xi|=1$.
    Then the same system has the solution
    \begin{align}
        \widehat B_i=\xi B_i,\qquad
        \widehat Y_j=\xi^{-j}Y_j,\qquad
        \widehat Z_j=\xi^{j+1}Z_j.
    \end{align}
    If $\xi^N=1$, the actual length-$N$ vectors of $B$ and
    $\widehat B$ are equal.
\end{lemma}
\begin{proof}
    The scalar product at a site is $\xi^{-j}\xi^{j+1}=\xi$,
    whereas the scalar product at an internal inverse equation is
    $\xi^{-(j+1)}\xi^{j+1}=1$. Unit modulus preserves row
    normalization. A length-$N$ periodic trace is multiplied by
    $\xi^N$, proving the vector equality.
\end{proof}

\begin{lemma}[Corrected converse for an actual same-state solution]
    \label{lem:quadratic_same_state_unitary_phase}
    \lean{MPSTensor.QuadraticReconstructionSolution.exists_unitary_phase_of_same_state,
        MPSTensor.QuadraticReconstructionSolution.unitGaugePhaseEquiv_of_same_state}
    \notready
    \uses{def:pgvwc07_quadratic_system, thm:pgvwc07_fixed_length_intertwiner}
    Let $A$ be a translation-invariant canonical tensor satisfying
    condition C1 at length $L_0$, and let $B$ be the tensor of a solution
    over a supplied OBC window. Assume
    \begin{align}
        N>2L_0+D^4,\qquad V^{(N)}(A)=V^{(N)}(B)\ne0.
    \end{align}
    Then $A$ is row normalized and there are a unitary $U$ and a scalar
    $\xi$ with $|\xi|=1$, $\xi^N=1$ and
    \begin{align}
        A_i=\xi U B_i U^\dagger
    \end{align}
    for every letter. No injectivity or canonical-form hypothesis on
    the solution tensor $B$ is added. The size condition provides both
    endpoint cuts through the actual fixed-length state equality;
    these cuts are not inferred from the internal quadratic equations.
\end{lemma}
\begin{proof}
    Condition C1 reduces the canonical tensor to one positive weighted
    block, with row sum $\lambda^2I$. The solution equation gives
    the row sum $I$ for $B$. The corrected weighted finite-length
    uniqueness theorem gives $\lambda=1$ and the displayed unitary
    conjugation and phase. Its complete-cut argument constructs all
    $D^4+1$ intertwiners, including the two endpoint cuts. Equality of
    the actual nonzero vectors forces $\xi^N=1$.
\end{proof}

\begin{lemma}[An actual rephased solution with exact unitary gauge]
    \label{lem:quadratic_same_state_exact_rephased_solution}
    \lean{MPSTensor.QuadraticReconstructionSolution.exists_rephased_solution_exact_gauge_of_same_state}
    \notready
    \uses{lem:quadratic_solution_rephasing, lem:quadratic_same_state_unitary_phase}
    Under the preceding hypotheses, the explicitly rephased solution
    over the same window has tensor $\widehat B_i=\xi B_i$, preserves
    the actual length-$N$ state, and satisfies
    \begin{align}
        A_i=U\widehat B_iU^\dagger.
    \end{align}
    In particular $\widehat B$ is exactly gauge equivalent to $A$.
    This is the corrected same-state converse; it does not assert that
    an arbitrary algebraic solution represents the prescribed state,
    or construct the solution from a unique OBC canonical window.
    The remaining source scope is recorded in
    \cite{gap:pgvwc07_quadratic_reconstruction_phase} and
    \cite{gap:pgvwc07_intertwiner_chain_off_by_one}.
\end{lemma}
\begin{proof}
    Apply the two preceding lemmas with the derived root of unity.
    A unitary matrix is a unit of the matrix ring, with its adjoint
    as inverse, so the same matrix gives the exact gauge equivalence.
\end{proof}

\subsection{Actual cut ranks and physical window gauges}
\label{sec:quadratic_physical_cut_window}

\begin{lemma}[Exact rank and basis of an injective physical cut]
    \label{lem:quadratic_injective_physical_cut}
    \lean{MPSTensor.cutMatrix_mpv_eq_cutVector_mul,
        MPSTensor.exists_rightInverse_rightCutVector,
        MPSTensor.span_cutMatrix_mpv_cols_eq,
        MPSTensor.periodicCutBasis, MPSTensor.periodicCutBasis_apply,
        MPSTensor.cutRank_mpv_eq_of_isNBlkInjective,
        MPSTensor.isNBlkInjective_of_le_of_canonical,
        MPSTensor.cutRank_mpv_eq_of_canonical,
        MPSTensor.cutCoefficientRank_mpv_eq_of_canonical,
        MPSTensor.quadraticWindow_cutCoefficientRank_eq}
    \notready
    \uses{thm:pgvwc07_fixed_length_intertwiner}
    Let $A$ be a canonical tensor of positive bond dimension satisfying
    condition C1 at length $L_0$. Every physical cut at $k$ with
    $L_0\leq k\leq N-L_0$ has Schmidt rank exactly $D^2$.
    Its prefix column space has the actual left virtual cut vectors
    as a basis. In particular all $D^4+1$ cuts of the quadratic
    reconstruction window have rank $D^2$ when $N>2L_0+D^4$.
\end{lemma}
\begin{proof}
    The coefficient matrix factors as $FG$, with columns of $F$ the
    left virtual vectors and rows of $G$ the right virtual vectors.
    Suffix injectivity gives a right inverse of $G$. Therefore the
    column spaces of $FG$ and $F$ agree. Prefix injectivity makes the
    $D^2$ columns of $F$ independent. A positive canonical block weight
    is removed internally to propagate C1 to every longer length, and
    then restored to retain the original tensor.
\end{proof}

\begin{lemma}[Physical slicing in the virtual cut bases]
    \label{lem:quadratic_virtual_cut_slicing}
    \lean{MPSTensor.physicalCutCoordinates,
        MPSTensor.physicalCutCoordinates_apply,
        MPSTensor.physicalPeriodicCutBasis,
        MPSTensor.physicalPeriodicCutBasis_apply,
        MPSTensor.toMatrix_cutSlice_physicalPeriodicCutBasis}
    \notready
    \uses{lem:quadratic_injective_physical_cut}
    Transfer the derived cut bases to the original physical-site
    coordinates. Fixing the next physical letter $i$ has matrix
    $A_i\otimes I_D$ between the two adjacent virtual cut bases.
\end{lemma}
\begin{proof}
    The configuration bijections identify the numbered and original
    physical column spaces. Appending the letter $i$ multiplies the
    virtual product by $A_i$ and leaves the second virtual index fixed.
\end{proof}

\begin{lemma}[A quadratic solution from the actual OBC window]
    \label{lem:quadratic_actual_physical_window_solution}
    \lean{MPSTensor.quadraticPhysicalCutWindow,
        MPSTensor.exists_physicalCutWindow_gauges,
        MPSTensor.exists_quadraticSolution_physicalCutWindow,
        MPSTensor.quadraticWindowCutBases,
        MPSTensor.exists_quadraticSolution_cutSiteWindow,
        MPSTensor.quadraticPhysicalCutWindow_cutSiteMatrix}
    \notready
    \uses{def:pgvwc07_quadratic_system, lem:quadratic_virtual_cut_slicing}
    Let $A$ be a canonical condition-C1 tensor as above. Choose any
    bases of its actual physical prefix column spaces. On an interior
    window, the physical slicing matrices in these bases admit a
    solution of the quadratic system. The window matrices are exactly
    the minimal open-boundary site matrices after the proved
    $D^2$-dimensional coordinate reindexing. No window gauges or
    quadratic solution are assumed.
\end{lemma}
\begin{proof}
    Let $C_j$ be the chosen physical cut bases and $Q_j$ the derived
    virtual bases. The matrices of the identity from $C_j$ to $Q_j$
    and from $Q_{j+1}$ to $C_{j+1}$ supply the two gauges. Composition
    of these identity maps gives each internal inverse relation and
    converts physical slicing to $A_i\otimes I_D$. If the canonical
    tensor has weight $\lambda>0$, rescale the two gauges by
    $\lambda^j$ and $\lambda^{-(j+1)}$. The resulting tensor is
    $\lambda^{-1}A$, which is row normalized. The physical window
    and the prescribed state used to form its column spaces are
    retained throughout this construction.
\end{proof}
